%% ============================================================
%%  Data in Brief – Article Outline
%%  Dataset: RNPDNO (Mexico, 2015–Oct 2025)
%%  Template: elsarticle, 3p option
%% ============================================================
\documentclass[3p]{elsarticle}

\usepackage{booktabs}
\usepackage{lineno}
\usepackage{hyperref}
\usepackage{array}
\usepackage{longtable}
\usepackage{caption}
\usepackage{microtype}
\usepackage{amsmath}
\usepackage{amssymb}
\usepackage{graphicx}
\usepackage{placeins}

\journal{Data in Brief}
\hypersetup{
  pdfauthor={Marco A. Escobar, Gerardo Reyes-Guzman, Abraham Sanchez Ruiz},
  pdftitle={A municipality-level, monthly dataset of disappearance cases and outcomes in Mexico from the National Registry of Disappeared Persons, 2015-2025},
  pdfkeywords={missing persons, personas desaparecidas, RNPDNO, administrative records, spatiotemporal analysis, human rights, open data},
  pdfsubject={Data in Brief submission}
}
%% ============================================================
\begin{document}

%\linenumbers
\begin{frontmatter}

\title{A municipality-level, monthly dataset of disappearance cases and outcomes
in Mexico from the National Registry of Disappeared Persons, 2015--2025}

\author[inst1]{Marco A. Escobar}
\author[inst2]{Gerardo Reyes-Guzman}
\author[inst3]{Abraham Sanchez Ruiz\corref{cor1}}
\cortext[cor1]{Corresponding author.}
\ead{asanchezr@lasallebajio.edu.mx}

\address[inst1]{Facultad de Ingenierías y Tecnologías, Universidad La Salle Bajío,
  León, Guanajuato 37150, Mexico}
\address[inst2]{Facultad de Economía, Benemérita Universidad Autónoma de Puebla,
  Puebla, Puebla 72570, Mexico}
\address[inst3]{Facultad de Estudios Superiores, Universidad La Salle Bajío,
  Salamanca, Guanajuato 37000, Mexico}


%% ----------------------------------------------------------
%%  ABSTRACT  (248 words; DIB limit = 250)
%% ----------------------------------------------------------
\begin{abstract}

Data were obtained from the publicly accessible sociodemographic
dashboard of Mexico's National Search Commission
(Comisi\'{o}n Nacional de B\'{u}squeda, CNB) through a Python-based
web scraper that systematically queried the National Registry of
Disappeared and Unlocated Persons (Registro Nacional de Personas
Desaparecidas y No Localizadas, RNPDNO) portal.
The scraper iterated over all combinations of state,
victim-status category, and monthly date ranges, retrieving
sex-disaggregated counts from HTML tables embedded in JSON API
responses.
Raw records were processed through a reproducible cleaning pipeline
to normalize administrative and geographic identifiers and to assign
standardized municipal geo-codes from Mexico's National Institute of
Statistics and Geography (Instituto Nacional de Estad\'{i}stica y
Geograf\'{i}a, INEGI).

The dataset is organized into four comma-separated value (CSV) files
corresponding to distinct victim-status categories:
(i)~all registered persons regardless of location outcome
(\texttt{status\_id}~=~0; 57{,}327 rows);
(ii)~persons disappeared and not yet located
(\texttt{status\_id}~=~7; 34{,}640 rows);
(iii)~persons located alive
(\texttt{status\_id}~=~2; 37{,}358 rows);
and
(iv)~persons located deceased
(\texttt{status\_id}~=~3; 9{,}883 rows).
Each row represents a unique municipality--year--month cell.
All files share a twelve-column schema comprising composite and
disaggregated INEGI geo-codes, normalized state and municipal names,
sex-disaggregated counts, total persons, calendar year, calendar
month, and victim-status identifier.
Three sentinel geocodes flag records with unresolvable geographic information, accounting for 4{,}136 rows (12{,}178 persons), retained for transparency.
The complete processing pipeline, including web scraping, text normalization and manual override files, is released alongside the dataset in an open-source code repository.
\end{abstract}
\begin{keyword}
missing persons \sep
RNPDNO \sep administrative records \sep
spatiotemporal analysis \sep human rights \sep open data
\end{keyword}


\end{frontmatter}

%% ============================================================
%%  SPECIFICATIONS TABLE  (mandatory)
%% ============================================================
\clearpage

\section*{Specifications Table}

\begin{table}[h!]
\centering
\caption*{}
\renewcommand{\arraystretch}{1.4}
\begin{tabular}{>{\bfseries}p{0.32\linewidth} p{0.62\linewidth}}
\toprule
Subject & Social Sciences; Criminology \\
\midrule
Specific subject area &
  Spatiotemporal analysis of disappearances and missing
  persons in Mexico at the municipality level \\
\midrule
Type of data &
  Table (CSV format; four files); intermediate collection files
  (Parquet format, one per victim-status category) \\
\midrule
Data collection &
  Administrative records from Mexico's National Search Commission
  (Comisi\'{o}n Nacional de B\'{u}squeda, CNB) obtained via the publicly
  accessible sociodemographic interface of the National Registry
  of Disappeared and Unlocated Persons (Registro Nacional de
  Personas Desaparecidas y No Localizadas, RNPDNO).
  A Python script using \texttt{requests} and \texttt{BeautifulSoup4}
  systematically queried the portal by state, victim-status
  category, and monthly date ranges (2015--2025), then parsed
  municipality-level counts from HTML tables embedded in JSON
  API responses.
  Processed through a deterministic cleaning pipeline
  with two-pass geographic normalization against the
  National Institute of Statistics and Geography (Instituto
  Nacional de Estad\'{i}stica y Geograf\'{i}a, INEGI) municipal
  geo-code catalog. \\
\midrule
Data source location &
  Country: Mexico.
  Institution: National Search Commission (Comisi\'{o}n Nacional de B\'{u}squeda, CNB).
  Source URL:
  \url{https://versionpublicarnpdno.segob.gob.mx/Dashboard/Sociodemografico} \\
\midrule
Data accessibility &
  Repository name: Open Science Framework (OSF) \newline
  Data identification number: 5qmwr \newline
  Direct URL to data: \url{https://osf.io/5qmwr/} \newline
  Code repository: \url{https://github.com/marco-escobar/rnpdno-scraper} \\
\midrule
Related research article &
  None (direct submission) \\
\bottomrule
\end{tabular}
\end{table}
\clearpage

%% ============================================================
%%  VALUE OF THE DATA
%% ============================================================

\section*{Value of the Data}

\begin{itemize}

  \item The RNPDNO portal employs session-based authentication and CAPTCHA verification that preclude bulk download or direct programmatic access; no public API or downloadable dataset is provided. This dataset offers the first machine-readable, analysis-ready extraction of the complete 2015--2025 registry, enabling reproducible computational research at national scale previously infeasible for independent researchers \cite{mandolessi_2022}.
  \item Monthly temporal resolution across 132 periods
    (January 2015 through December 2025) supports sub-annual
    analyses, including event-study designs around policy
    interventions, seasonal decomposition, and lagged response
    modeling, that are not possible with the annual or cumulative
    snapshots used in prior municipality-level studies
    \cite{cadena_garrocho_2019, atuesta_2024}.

  \item The standardized five-digit INEGI geo-code (\texttt{cvegeo})
    enables direct spatial joins to official cartographic boundaries
    (EPSG:6372), population projections, and census microdata,
    supporting population-normalized rate computation and spatial
    cluster detection methods (LISA, Getis-Ord $G_i^*$) without
    intermediate reconciliation \cite{anselin_1995, inegi_ageeml}.

  \item The four-file structure by victim-status category
    (total registered, disappeared and not located, located alive,
    located deceased) permits outcome-transition analysis:
    researchers can compute location rates, resolution-time
    proxies, and spatial heterogeneity in case outcomes, analytic
    dimensions limited to single cities in existing work
    \cite{aguilar_velazquez_2025}.

  \item The open-source processing pipeline, including scraper,
    text normalizer, two-pass geographic join, and git-tracked
    manual override files, can be re-executed against future
    RNPDNO releases to extend temporal coverage; schema-validation
    modules detect source drift, supporting longitudinal research
    programs aligned with emerging data-based disappearance
    analysis frameworks \cite{ruiz_reyes_2024}.

\end{itemize}

\FloatBarrier
\clearpage

\section*{Data Description}

The dataset consists of four processed CSV files, one per RNPDNO
victim-status category, and four corresponding intermediate Parquet
files retained from the collection stage.

\subsection*{Processed CSV files}

All four processed files share the twelve-column schema described in
Table~\ref{tab:schema}.
Each row represents one (municipality, year, month, status) observation.

\begin{table}[ht]
  \centering
  \caption{Schema of the four processed RNPDNO CSV files.
    All four files (\texttt{rnpdno\_total.csv},
    \texttt{rnpdno\_disappeared\_not\_located.csv},
    \texttt{rnpdno\_located\_alive.csv},
    \texttt{rnpdno\_located\_dead.csv})
    share this schema.
    Geo-codes conform to the INEGI municipal geo-code catalog.
    Counts represent persons reported within the given
    municipality, month, and year.}
  \label{tab:schema}
  \renewcommand{\arraystretch}{1.2}
  \begin{tabular}{llp{0.45\linewidth}}
    \toprule
    Column & Type & Description \\
    \midrule
    \texttt{cvegeo}     & string  & Five-digit composite INEGI geo-code
                                    (\texttt{CVE\_ENT} + \texttt{CVE\_MUN});
                                    zero-padded.
                                    Sentinel values: \texttt{99998}
                                    (unknown state),
                                    \texttt{XX998} / \texttt{XX999}
                                    (municipality unspecified/unknown). \\
    \texttt{cve\_estado}& string  & Two-digit INEGI state code,
                                    zero-padded (01--32).
                                    Value \texttt{99} reserved for
                                    unknown state. \\
    \texttt{state}      & string  & Official state name (Spanish),
                                    normalized from the INEGI catalog. \\
    \texttt{cve\_mun}   & string  & Three-digit INEGI municipal code,
                                    zero-padded (001--570).
                                    Sentinel values: \texttt{998}
                                    (no municipal reference),
                                    \texttt{999} (municipality unknown). \\
    \texttt{municipality}& string & Official municipal name (Spanish),
                                    normalized from the INEGI catalog. \\
    \texttt{male}       & integer & Count of male persons. \\
    \texttt{female}     & integer & Count of female persons. \\
    \texttt{undefined}  & integer & Count of persons with indeterminate
                                    or unreported sex. \\
    \texttt{total}      & integer & Total persons
                                    (\texttt{male} + \texttt{female}
                                    + \texttt{undefined}). \\
    \texttt{year}       & integer & Calendar year (2015--2025). \\
    \texttt{month}      & integer & Calendar month (1--12). \\
    \texttt{status\_id} & integer & RNPDNO victim-status code:
                                    0~=~total, 7~=~disappeared and
                                    not located, 2~=~located alive,
                                    3~=~located deceased. \\
    \bottomrule
  \end{tabular}
\end{table}

\paragraph{\texttt{rnpdno\_total.csv}}
Contains 57{,}327 rows corresponding to all persons registered in
the RNPDNO, regardless of subsequent location status
(\texttt{status\_id}~=~0).
Temporal coverage spans January 2015 through December 2025
(up to 132 monthly periods per municipality).

\paragraph{\texttt{rnpdno\_disappeared\_not\_located.csv}}
Contains 34{,}640 rows for persons classified as disappeared and
not yet located (\texttt{status\_id}~=~7).
This file is the primary analytical target for
spatiotemporal characterization of active disappearance caseloads.

\paragraph{\texttt{rnpdno\_located\_alive.csv}}
Contains 37{,}358 rows for persons subsequently located alive
(\texttt{status\_id}~=~2).

\paragraph{\texttt{rnpdno\_located\_dead.csv}}
Contains 9{,}883 rows for persons located deceased
(\texttt{status\_id}~=~3).


\subsection*{Overview figures}

Figure~\ref{fig:trend} shows the monthly national count of persons
disappeared and not located (status\,=\,7) over the full 2015--2025
period, disaggregated by sex.
Figure~\ref{fig:states} shows the cumulative total for all statuses
(status\,=\,0) by state, ranked by total volume.

\begin{figure}[ht]
  \centering
  \includegraphics[width=\linewidth]{fig1_monthly_trend.pdf}
\caption{Monthly national count of persons disappeared and not located
    (\texttt{status\_id}\,=\,7), January 2015 -- December 2025,
    disaggregated by sex.
    Records with unknown state are excluded.}
  \label{fig:trend}
\end{figure}

\begin{figure}[ht]
  \centering
  \includegraphics[width=\linewidth]{fig2_state_distribution.pdf}
\caption{Cumulative total of registered persons across all outcome categories by state
    (\texttt{status\_id}\,=\,0), 2015--2025, disaggregated by sex.
    States ranked by total volume (ascending).
    Records with unknown state are excluded.}
  \label{fig:states}
\end{figure}

\subsection*{Sentinel geo-code records}

A subset of rows in each file carries sentinel geo-codes indicating
unresolvable geographic information in the source data
(Table~\ref{tab:sentinels}).
These rows are retained rather than discarded to preserve
the original aggregate counts; analysts performing
municipality-level spatial analyses should filter them as
appropriate.
\begin{table}[ht]
  \centering
  \caption{Sentinel geo-codes assigned to records with
    unresolvable geographic information, summed across all four
    processed CSV files. A given municipality-month cell may appear
    in multiple files if cases with different outcomes were registered
    there.}
  \label{tab:sentinels}
  \renewcommand{\arraystretch}{1.2}
  \begin{tabular}{llrr}
    \toprule
    \texttt{cvegeo} & Condition & Rows & Persons \\
    \midrule
    \texttt{99998} & Unknown state
                   & 198 & 408 \\
    \texttt{XX998} & No municipal reference
                   & 1{,}748 & 3{,}978 \\
    \texttt{XX999} & Municipality unknown
                   & 2{,}190 & 7{,}792 \\
    \midrule
    \textbf{Total} & & \textbf{4{,}136} & \textbf{12{,}178} \\
    \bottomrule
  \end{tabular}
\end{table}

\subsection*{Intermediate Parquet files}

Four intermediate Parquet files contain the consolidated raw
records after scraping but prior to geographic normalization.
These files retain the original column names and free-text
municipality strings from the RNPDNO portal and are provided
for provenance purposes.
Filenames reflect the original Spanish terminology used in the source system.

\begin{description}
  \item[status 0:] \texttt{rnpdno\_total\_2015\_2025.parquet}
  \item[status 2:] \texttt{rnpdno\_localizadas\_con\_vida\_2015\_2025.parquet}
  \item[status 3:] \texttt{rnpdno\_localizadas\_sin\_vida\_2015\_2025.parquet}
  \item[status 7:] \texttt{rnpdno\_desaparecidas\_no\_localizadas\_2015\_2025.parquet}
\end{description}


\subsection*{Quality report}

The file \texttt{3\_documentation/quality\_report.csv} contains
per-dataset quality metrics generated at the end of each pipeline
run, including row counts, sex-total consistency checks, and
sentinel-record counts.
As of the reference extraction (March 2026), all four processed files
pass schema validation; zero sex-total inconsistencies are present;
and exact duplicate rows originating from source-data artefacts
have been removed.
Sentinel-coded records (Table~\ref{tab:sentinels}) reflect
unresolvable geographic information in the original RNPDNO source
data, not pipeline matching failures.
\clearpage
%% ============================================================
%%  EXPERIMENTAL DESIGN, MATERIALS, AND METHODS
%% ============================================================
\section*{Experimental Design, Materials, and Methods}

\subsection*{Data source and access}

Data were obtained from the publicly accessible sociodemographic
query interface of the RNPDNO, maintained by the Comisión Nacional
de Búsqueda (CNB), available at:
\url{https://versionpublicarnpdno.segob.gob.mx/SocioDemografico/TablaDetalle}.
The interface presents HTML tables aggregated by Mexican state
(INEGI codes 01--32), year (2015--2025), and calendar month (1--12),
for four victim-status categories
(\texttt{id\_estatus\_victima}: 0, 2, 3, 7).
No bulk-download functionality is provided; access requires
iterative HTTP requests with a valid session cookie.

\subsection*{Web scraping}

The script \texttt{scripts/scrapers/scrape\_rnpdno\_single\_status.py} (Python~3.10)
implements a single-status scraper that iterates over all
combinations of state $\times$ year $\times$ month for one
status value, using the \texttt{requests} library for HTTP
communication and \texttt{BeautifulSoup4} for parsing HTML tables
embedded in JSON API responses.
Requests are serialized with a 0.5~s inter-state and 1.0~s
inter-month delay to avoid server overload; a retry mechanism
with exponential backoff (up to four retries, 45~s timeout)
handles transient failures.
Four independent invocations produce one Parquet file per
status value, consolidated from individual CSV exports per
state--year--month triplet.
State~33 (\emph{Se desconoce}) is included in the iteration
to capture records with unresolvable state-level geography,
which are subsequently assigned sentinel geo-code
\texttt{cvegeo}~=~\texttt{99998}.

\subsection*{Geographic reference construction}

The INEGI \emph{Archivo General de la Entidad Marco de Localidades}
(AGEEML), obtained from \cite{inegi_ageeml}, provides the canonical
municipality-level geographic reference.
The script \texttt{src/build\_reference.py} reads the raw AGEEML
file (ISO-8859-1 encoding, 2{,}478 rows), zero-pads geo-codes to
canonical widths (\texttt{CVEGEO}: 5 digits; \texttt{CVE\_ENT}:
2 digits; \texttt{CVE\_MUN}: 3 digits), and derives two normalized
name variants using utilities in \texttt{src/normalize.py}:

\begin{description}
  \item[Standard normalization]
    Lower-case; strip Unicode accents via NFD decomposition;
    remove punctuation; collapse whitespace.
    Formally:
    \[
      f_{\text{std}}(s) = \text{STRIP}\!\left(
  \text{COLLAPSE}\!\left(
    \text{RM\_PUNCT}\!\left(
      \text{NFD\_STRIP}(s.\text{lower}())
    \right)
  \right)
\right)
    \]

  \item[Aggressive normalization]
    Applies $f_{\text{std}}$ then expands a curated abbreviation
    dictionary (e.g., \texttt{Dr.}$\to$\texttt{doctor},
    \texttt{Mpio.}$\to$\texttt{municipio}), removes common
    prepositions (\emph{de}, \emph{del}, \emph{los}, \emph{las}),
    and collapses repeated whitespace.
\end{description}

\subsection*{Pre-join state corrections (B1 heuristic)}

The script \texttt{src/build\_state\_corrections\_b1.py}
identifies municipality names that are unique nationwide in
AGEEML (i.e., appear in exactly one state).
For raw RNPDNO records whose reported state does not match the
unique AGEEML state for their municipality name, the state
identifier is corrected prior to the geographic join.
This procedure produced 25 validated corrections
(rule: \texttt{B1\_UNIQUE\_NATIONWIDE\_EXACT}),
stored in \texttt{data/manual/geo\_state\_corrections.csv}.

\subsection*{Two-pass geographic join}

The main pipeline (\texttt{src/clean\_pipeline.py}) matches each
raw record to an AGEEML entry via a two-pass procedure:

\begin{enumerate}
  \item \textbf{Pass 1 – standard join.}
    Inner join on $(\texttt{cve\_estado}, f_{\text{std}}(\texttt{municipio}))$.

  \item \textbf{Pass 2 – aggressive join (unambiguous only).}
    For records unmatched in Pass~1, inner join on
    $(\texttt{cve\_estado},$ $f_{\text{agg}}(\texttt{municipio}))$,
    restricted to keys that map to exactly one AGEEML entry
    (i.e., no ambiguity).
\end{enumerate}

Records unmatched after both passes are written to
\texttt{data/manual/unmatched\_rnpdno\_*.csv} for manual review.

\subsection*{Manual overrides}

A git-tracked file \texttt{data/manual/geo\_overrides.csv} (35 rows)
maps $(raw\_id\_estado, raw\_municipio) \to \texttt{cvegeo}$
for cases requiring human disambiguation
(e.g., politically renamed municipalities, cross-state
boundary ambiguities).
Overrides are applied after both join passes and before
sentinel assignment.

\subsection*{Sentinel geo-code assignment}

Records that remain unresolvable after overrides are assigned
sentinel geo-codes according to the logic in
\texttt{data/manual/special\_geo\_cases.csv}:

\begin{description}
  \item[\texttt{cvegeo = 99998}]
    Assigned to records with \texttt{id\_estado}~=~33
    (unknown state) in the source data.

  \item[\texttt{cve\_mun = 998}]
    Source municipality string indicates no municipal reference.

  \item[\texttt{cve\_mun = 999}]
    Source municipality string indicates municipality unknown.
\end{description}


\subsection*{Quality checks}

The pipeline executes the following automated quality checks
before writing output:

\begin{itemize}
  \item Schema validation: all twelve columns present with
    expected types.
  \item Duplicate detection: exact-row duplicates originating
    from source-data artefacts are identified and removed.
  \item Sex-total consistency:
    $\texttt{total} = \texttt{male} + \texttt{female}
    + \texttt{undefined}$ verified for every row.
  \item Temporal coverage: year range [2015,~2025] and month
    range [1,~12] verified.
  \item Cross-status consistency: the aggregate person-count
    in the status-0 (total) file is compared against the sum
    of the three sub-status files
    (\texttt{status\_id}~$\in\{2,3,7\}$).
    As of the reference scrape, the total file records
    268{,}900 persons while the sum of the three sub-status
    files yields 268{,}973 persons---the total file contains
    73 fewer persons than the sub-status sum ($0.027\%$),
    distributed across 55{,}524
    (municipality, year, month) cells and confined to
    deviations of $\pm1$ or $\pm2$ persons per cell.
    This near-exact agreement confirms internal consistency
    of the source system; the residual discrepancy is
    attributed to asynchronous updates in the RNPDNO portal.
\end{itemize}

Results are written to \texttt{3\_documentation/quality\_report.csv}.


\subsection*{Output assembly}

Two municipality names could not be resolved to unique INEGI
geo-codes across either join pass or the manual override file,
and are therefore retained in the processed files with a null
\texttt{cvegeo} value rather than discarded, preserving the
full aggregate person-counts.

\textbf{PUEBLO NUEVO} appears in the RNPDNO source under
three states: Durango (\texttt{id\_estado}~=~10),
Guanajuato (\texttt{id\_estado}~=~11), and Sinaloa
(\texttt{id\_estado}~=~25).
AGEEML contains canonical entries for Durango
(\texttt{cvegeo}~=~\texttt{10023}) and Guanajuato
(\texttt{cvegeo}~=~\texttt{11024}), but no municipality named
Pueblo Nuevo exists in Sinaloa, suggesting a
source-data error for the Sinaloa occurrence.
All three (estado, \texttt{PUEBLO NUEVO}) combinations were
flagged for manual review; the geo\_overrides file was not
completed for these entries and they remain unresolved.

\textbf{NOGALES} appears under Sinaloa
(\texttt{id\_estado}~=~25), Sonora
(\texttt{id\_estado}~=~26), and Veracruz
(\texttt{id\_estado}~=~30).
AGEEML contains canonical entries for Sonora
(\texttt{cvegeo}~=~\texttt{26043}) and Veracruz
(\texttt{cvegeo}~=~\texttt{30115}), but no municipality named
Nogales exists in Sinaloa.
Similarly, all three occurrences were flagged for manual
review and the overrides file was left incomplete.

Together, these six (estado, municipio) combinations account
for 237 rows and 1{,}361 persons in the status-0 file; the
corresponding unresolved rows in the three sub-status files
total 380 rows (176, 144, and 60 for statuses 7, 2, and 3
respectively).
Analysts performing municipality-level spatial analyses should
exclude rows with null \texttt{cvegeo} as appropriate.

The four processed CSV files are written to
\texttt{data/processed/}.

\subsection*{Software environment}

The pipeline requires Python~3.10 with the following
packages: \texttt{pandas}~$\geq$2.0, \texttt{pyarrow}~$\geq$14.0,
\texttt{requests}~$\geq$2.31, \texttt{beautifulsoup4}~$\geq$4.12,
\texttt{python-dotenv}~$\geq$1.0.
Code quality is enforced with \texttt{ruff} (line-length~=~99).
Environment reproducibility is managed via a Conda environment
file (\texttt{environment.yml}) and a DVC pipeline configuration.
All code and manual override files are version-controlled in a
public git repository \cite{code_repo}.

%% ============================================================
%%  ETHICS STATEMENT
%% ============================================================
\section*{Ethics Statement}

The RNPDNO dataset is derived exclusively from publicly accessible
administrative records published by the Mexican federal government
through the Comisión Nacional de Búsqueda.
All data are presented in aggregated, municipality-level form;
no individual-level identifiers (names, dates of birth, case
numbers, or biometric data) are present in any of the processed
files.
Accordingly, no Institutional Review Board (IRB) approval or
ethics committee clearance is required for secondary analysis
of these public administrative data.
The collection and processing activities described here are
consistent with applicable data-use conditions stated by the CNB
for publicly disseminated RNPDNO statistics.

%% ============================================================
%%  CREDIT AUTHOR STATEMENT
%% ============================================================
\section*{CRediT Author Statement}

\textbf{Marco A. Escobar:} Conceptualization, Data curation,
Formal analysis, Software, Writing -- original draft.
\textbf{Gerardo Reyes-Guzman:} Validation,
Writing -- review \& editing, Supervision.
\textbf{Abraham Sanchez Ruiz:} Validation,
Writing -- review \& editing, Supervision.


%% ============================================================
%%  DECLARATION OF COMPETING INTERESTS
%% ============================================================
\section*{Declaration of Competing Interests}

The authors declare that they have no known competing financial
interests or personal relationships that could have influenced
the work reported in this data article.

%% ============================================================
%%  ACKNOWLEDGMENTS / FUNDING
%% ============================================================
\section*{Acknowledgments}

\textbf{Funding:} This research received financial support from
Universidad La Salle Bajío (Grant No. ITC-2025-18) for data acquisition and processing.

The authors acknowledge the Comisión Nacional de Búsqueda for
publicly releasing the RNPDNO aggregate statistics.

\section*{Declaration of generative AI and AI-assisted technologies in the manuscript preparation process}

During the preparation of this work, the authors used Claude
(Anthropic) to assist with the refinement of manuscript text for
clarity and conciseness.
After using this tool, the authors reviewed and edited the content
as needed and take full responsibility for the content of the
publication.



%% ============================================================
%%  REFERENCES
%% ============================================================
\bibliographystyle{elsarticle-num}
\bibliography{references}

\end{document}
